Sea $f(x, y, z) = x^2 + y^2 + z^2$ y sea $\mathbf{r} = \langle x, y, z \rangle$ con $r = \|\mathbf{r}\|$.
\begin{enumerate}
    \item Muestre que $\nabla(r^n) = n r^{n-2} \mathbf{r}$ para cualquier entero $n \geq 1$.
    \item Deduzca que $\nabla(1/r) = -\mathbf{r}/r^3$ (campo gravitacional de una masa puntual, salvo constante).
    \item Una función $f(\mathbf{r})$ cuyo gradiente apunta en la dirección radial $\mathbf{r}/r$ se llama función central. Demuestre que toda función de la forma $f = \phi(r)$, con $\phi$ diferenciable, es una función central.
\end{enumerate}
